\section{从动态规划到 Q 学习}
\label{sec:16-Control-and-Reinforcement-Learning/02-Reinforcement-Learning/04}

Bellman 最优性方程已经把目标写清楚了：

\[
Q^*(s,a)
=
\E\left[
R_{t+1}
+\gamma\max_{a'}Q^*(S_{t+1},a')
\mid S_t=s,A_t=a
\right].
\]

问题是右端的期望怎么算。动态规划的做法是打开转移模型：对每个可能的下一状态 \(s'\)，查出 \(P(s'\mid s,a)\)，加权求和。真实机器人通常没有这张表——做了动作，传感器只返回\emph{一个}下一状态和\emph{一份}奖励，没有概率分布。

能不能用这一个随机样本，替代对所有未来的加权求和？Q 学习的回答是：可以，但只能逐步修正。每次交互都给你一份带噪声的 Bellman 目标；学习率把无数次局部修正平均起来，最终逼近同一个不动点。

\subsection{从 Bellman 残差到时序差分误差}

在时刻 \(t\)，你观察到了一个转移

\[
(S_t,A_t,R_{t+1},S_{t+1}),
\]

用当前的 \(Q\) 估计构造一个目标值

\[
Y_t
=
R_{t+1}
+\gamma\max_{a'}Q_t(S_{t+1},a').
\]

目标与当前预测之间的差距

\[
\delta_t
=Y_t-Q_t(S_t,A_t)
\]

叫时序差分（TD）误差。Q 学习只更新刚刚访问过的那个状态动作对：

\[
Q_{t+1}(S_t,A_t)
=
Q_t(S_t,A_t)
+\alpha_t\delta_t,
\]

其他表项保持不变。\(\delta_t>0\) 说明这次结果比预期好，往上调；\(\delta_t<0\) 说明比预期差，往下调。

这个更新在形式上很像随机梯度下降。但有一个关键差别：目标 \(Y_t\) 里用了当前的 \(Q_t(S_{t+1},a')\)——你用的是你自己的估计去生成新的估计。这叫 bootstrap：用已有知识构造学习目标，再用目标去修正知识。

在 \((S_t,A_t)\) 上取条件期望，\(Y_t\) 正好是 Bellman 最优性算子作用在 \(Q_t\) 上的结果。单次样本有噪声，但没有系统性的偏离。Q 学习可以看成异步、随机的值迭代——每一步只更新一个表项，用采样代替期望。

\subsection{离策略：收集数据的手和做决策的手可以不同}

收集数据的策略叫行为策略，记为 \(\mu\)。它可以为了探索而随机选动作——实际上\emph{必须}随机，否则很多动作永远没数据。但更新目标里用的是

\[
\max_{a'}Q_t(S_{t+1},a'),
\]

它假设了下一步会按当前贪心策略（取最大值）行动。数据是 \(\mu\) 采集的，目标是按贪心策略算的——两只手不是同一个策略。所以 Q 学习是离策略（off-policy）方法。

和 SARSA 比着看最清楚。SARSA 的更新用的是实际执行的下一动作 \(A_{t+1}\)：

\[
Q(S_t,A_t)
\leftarrow
Q(S_t,A_t)
+\alpha_t
\left[
R_{t+1}+\gamma Q(S_{t+1},A_{t+1})-Q(S_t,A_t)
\right].
\]

它评价的是\emph{包含探索行为在内的当前策略}，属于同策略（on-policy）学习。为什么这个差别要紧？如果一个探索动作有真实风险——比如机器人试着走一条不确定的窄道——Q 学习的目标假设未来会贪心（不走那条窄道），SARSA 会把未来仍可能随机探索造成的风险纳入价值。离策略不是天然更好：它只是允许``用一种策略收数据，同时学习另一种策略''。

\subsection{Expected Sarsa：先平均掉下一动作的随机性}

SARSA 的目标里还要随机抽一次 \(A_{t+1}\)。如果当前策略 \(\pi(\cdot\mid S_{t+1})\) 已知，可以在看到下一状态后直接对动作求条件期望：

\[
Y_t^{\mathrm{ES}}
=R_{t+1}+\gamma
\sum_{a'}\pi(a'\mid S_{t+1})Q_t(S_{t+1},a').
\]

这就是 Expected Sarsa。它和普通 SARSA 在给定 \(S_{t+1}\) 时有相同的期望，但消除了``下一动作恰好抽到哪一个''带来的那部分方差。动作数不多时，一次精确求和的代价换方差降低，通常划算；动作空间巨大或连续时，精确求和本身就成了新的计算难题。

三种目标的差异可以精确标定。SARSA 用实际抽到的动作；Expected Sarsa 对当前策略的动作分布求平均；Q 学习用最大值。前两者评价给定策略，最后一个直接逼近最优性算子。把平均换成最大值不只改变方差，也把要解的 Bellman 方程本身换掉了。比方外卖骑手在雨天选路线：Expected Sarsa 会把当前策略偶尔走拥堵小路的概率一起计入长期耗时；Q 学习的目标却假定下一路口以后总会走当前看来最快的方向。

\subsection{探索：没试过的动作没有数据}

如果从一开始就总选当前 \(\argmax_a Q(s,a)\)，初始估计的偶然偏差可能让某些动作永远不被尝试，错误便无法被纠正。最简单的 \(\varepsilon\)-greedy 策略：以概率 \(1-\varepsilon\) 选当前贪心动作，以概率 \(\varepsilon\) 在全部动作中均匀随机探索。

这个困境在日常生活中常见。常去一家面馆，永远只点最熟悉的那碗面。熟悉菜品的评价来自许多次品尝，新菜品的评价却只是猜测；每次选当前评分最高的，新菜品永远没有数据，``它不如旧菜''的判断从未被真正检验。偶尔换一种可能损失一顿饭的满意度，却换来以后选择所需的信息。在一般强化学习里，尝试一个动作还会把你带到新的后继状态，信息收益和后续代价都会继续传播。

探索和利用形成了真实的冲突。探索牺牲眼前回报换取信息，利用把已有知识变成回报。训练用的模拟器里可以大量探索；医疗、金融和实体机器人里，坏动作的代价不能用``以后会学会''来抵消，需要安全约束、离线数据、保守估计或人工规则层层把关。MDP 的奖励如果没有把安全红线写对，单纯加大碰撞罚分也未必足以阻止罕见灾难。

理论的收敛通常要求每个状态动作对被访问无穷多次，同时策略最终趋于贪心。这可以由逐渐衰减但不太快的 \(\varepsilon_t\) 实现——比如 \(\varepsilon_t = 1/t\) 级别。固定较大的探索率不会收敛到纯贪心行为；衰减太快又可能永远漏掉某些关键动作。

\subsection{表格情形什么时候真的收敛}

对有限 MDP，奖励有界、\(0\le\gamma<1\)，若每个状态动作对被访问无穷多次，且学习率对每一表项满足 Robbins--Monro 条件

\[
\sum_t\alpha_t=\infty,
\qquad
\sum_t\alpha_t^2<\infty,
\]

则表格 Q 学习在适当条件下几乎必然收敛到 \(Q^*\)。第一条保证新数据永远有累积影响力——旧数据不能把新数据完全淹没；第二条保证噪声的累积方差有限——单次随机波动最终被平均掉。典型选择是某表项第 \(n\) 次被访问时取 \(\alpha_n=1/n\)。

这两个条件也解释了常见的工程现象。固定学习率不会严格收敛到一个常数，而会持续追踪并在真值附近波动——在非平稳环境中这反而有用，因为旧数据必须逐渐被遗忘。理论里衰减步长适合固定世界，工程里固定步长适合移动世界，选择不同是因为面对的问题不同。

收敛定理也只针对表格情形：每个 \((s,a)\) 有一个独立可修改的数值。它没有告诉你有限数据时误差有多大，没有告诉你探索是否安全，也没有说把表格换成神经网络后同样的更新规则仍然收敛。

\subsection{一个数字例子：把更新走一遍}

假设一个状态动作对的当前估计 \(Q(s,a)=4\)。你执行了动作，收到即时奖励 \(r=1\)，观察到下一状态。在下一状态查表，所有动作中最大估值为 \(5\)。取 \(\gamma=0.9\)、\(\alpha=0.2\)。

目标是

\[
y=1+0.9\times5=5.5,
\]

TD 误差是 \(\delta=5.5-4=1.5\)，更新后

\[
Q_{\mathrm{new}}(s,a)
=4+0.2\times1.5
=4.3.
\]

这个例子里有两个容易踩的坑。第一个：不要把 \(5.5\) 当成观测到的真实总回报。观测到的是即时奖励 \(1\) 和下一状态本身；后面的 \(4.5\)（\(0.9\times5\)）来自你当前模型自己的估计——bootstrap 把现有知识掺进了目标。这提高了样本效率（一个样本同时更新了对当前和未来的认识），也让误差可以自我传导：如果下一状态被严重高估，本状态会沿着 Bellman 备份继承那份乐观。

第二个坑更隐蔽：取最大值操作。多个动作的估计都带噪声时，\(\max\) 更容易挑中正噪声偏大的那个，使 Q 值系统性偏高。这是选择偏差——同一份噪声既负责选动作又负责打分。Double Q-learning 把这两件事分开：一套估计用来选动作，另一套用来评价动作，降低``既当裁判又当运动员''造成的膨胀。

\subsection{表格结论不能直接搬到函数空间}

表格 Q 学习的收敛性依赖一个关键前提：每个状态动作对拥有可单独修改、互不干扰的表项。你把表换成参数函数 \(Q(s,a;\theta)\) 后，一次参数更新会同时改变许多状态的估计。如果目标又来自当前估计（bootstrap），误差就可能沿着共享参数反复放大。行为策略与目标策略不一致时，训练数据强调的状态还可能与待评价策略真正访问的状态不同——数据分布不匹配。

后面四节把这条边界逐层拆开：先在线性函数空间中说明 TD 实际求解的是投影 Bellman 方程；再用一个可计算的二维例子展示 Deadly Triad（函数逼近、bootstrap、离策略训练）怎样让误差放大；随后解释 DQN 的经验回放和目标网络分别减慢哪一类变化；最后用重要性采样精确区分 on-policy、off-policy、online 与 offline。这些都不是 Q 学习公式里的实现细节——它们是从表格算法走到大状态空间时\emph{必须增加}的条件。

\subsection{同一个不动点，不同的信息条件}

把两个单元串起来，结构已经很清楚了：

\begin{itemize}
\tightlist
\item
  HJB 用已知连续动力学写出无穷小时间片上的 Bellman 方程；
\item
  离散动态规划用已知转移概率做精确的 Bellman 备份——对全部下一状态加权求和；
\item
  Q 学习用单次样本转移做随机 Bellman 备份——一次只看一个实际到达的下一状态；
\item
  函数逼近再把装不下的值函数表压进有限参数。
\end{itemize}

每向现实靠近一步，就少一项信息、多一项近似，也多一种可能失效的方式。强化学习并不保证``只要试得够久就会变聪明''：状态要近似 Markov，奖励要表达真实目标，探索要覆盖到关键动作，数据分布要撑得起待评估的策略，函数逼近和 bootstrap 还要保持稳定。Bellman 方程给出了方向，可信的学习结果取决于这些前提是否真的成立。
